\documentclass{beamer}
%%%%%%%%%%%%%%%%%%%%%%%%%%%%%%%%%%
% CONFIGURACIÓN DEL TEMA
%%%%%%%%%%%%%%%%%%%%%%%%%%%%%%%%%%
\usetheme{metropolis}
%\usetheme[progressbar=frametitle]{metropolis}
%\definecolor{greenbeamer}{rgb}{0.125,0.5,0.25}
%\setbeamertemplate{frame numbering}[fraction] % para que diga las 
%\useoutertheme{metropolis}
%\useinnertheme{metropolis}
%\usefonttheme{metropolis}
%\usecolortheme{spruce}
%\setbeamercolor{background canvas}{bg=white}
%\usecolortheme[named=greenbeamer]{structure}


\usepackage[spanish, es-noquoting]{babel} %necesario para el idioma
%\usepackage{amsmath}
\usepackage{mathtools}
\usepackage{tikz} % necesario para tikz
\usepackage{chemmacros}
%\usepackage{chemformula} % necesario para fórmulas
\usepackage{enumitem} % necesario para personalizar las etiquetas de la enumeración
\usepackage{mycommands} % aquí tengo mis comandos
\usepackage{mycolours} % Aquí tengo mis colores.
\usepackage{mystyle_beamer_inorganica}

\usepackage{worldflags}

\setbeamercovered{transparent}


\definecolor{mDarkBrown}{HTML}{604c38}
\definecolor{mDarkTeal}{HTML}{23373b}
\definecolor{mLightBrown}{HTML}{EB811B}
\definecolor{mLightGreen}{HTML}{14B03D}

\newcommand*{\numberballmDarkTeal}{\protect\cnumero{\arabic*}{mDarkTeal}}
\newcommand*{\numberballmDarkBrown}{\protect\cnumero{\arabic*}{mDarkBrown}}
\newcommand*{\numberballmLightBrown}{\protect\cnumero{\arabic*}{mLightBrown}}
\newcommand*{\numberballmLightGreen}{\protect\cnumero{\arabic*}{mLightGreen}}

\newcommand{\examplebeamer}[1]{{\color{mLightGreen}\textbf{#1}}}
\newcommand{\alertbeamer}[1]{\alert{\textbf{#1}}}
\newcommand{\blockbeamer}[1]{{\color{mDarkTeal}\textbf{#1}}}

\newcommand{\numeromDarkTeal}[1]{\tikz[baseline=(myanchor.base)]
\node[minimum size=0.2cm,circle,
inner sep=1pt,draw, mDarkTeal,thick,fill=mDarkTeal] (myanchor)
{\color{white}\bfseries\fontsize{8}{8}#1};}



\newcommand{\ies}{IES Escultor José Luís Sánchez}

\newcommand{\eo}[3]{$\overset{^{\color{#1}#2}}{\text{\ch{#3}}}$}
\newcommand{\eob}[2]{$\overset{^{\color{blue}#1}}{\text{\ch{#2}}}$}
\newcommand{\eor}[2]{$\overset{^{\color{red}#1}}{\text{\ch{#2}}}$}

\newcommand{\molde}[2]{#1\text{ mol de }\ch{#2}}
\newcommand{\gde}[2]{#1\ \unit{\g} \text{ de }\ch{#2}}
\newcommand{\nline}[1][6pt]{\\[#1]}
\newcommand{\gasesideal}{P\ V = n\ R\ T}


\title{Formulación Química Inorgánica}
\subtitle{Normas IUPAC 2005}
\author{}
\institute{\textbf{Tipos de compuestos} Compuestos binarios: Óxidos, Hidruros, Peróxidos, Hidróxidos. Oxoácidos. Oxosales}
\date{}

\begin{document}


\begin{frame}
    \titlepage
\end{frame}

\begin{frame}
    \tableofcontents
\end{frame}


\section{Compuestos Binarios}

\subsection{Compuestos con oxígeno}

\begin{frame}[t]{Compuestos binarios con oxígeno: óxidos}\vspace{10pt}
    \begin{block}{\Large Características}
        \begin{enumerate}[label=\numberballmDarkTeal]
            \item Formados por \ch{O} y un elemento metálico o no metálico a excepción de los \alert{halógenos}.
            \item En este tipo de compuestos el oxígeno tiene un \azul{estado de oxidación ($-2$)}.
            \item Fórmula General: \ch{M_{x}O_{y}}
        \end{enumerate}
    \end{block}


    % \begin{alertblock}{\Large Nomenclatura}
    %     \vspace*{2mm}
    %     \begin{enumerate}[label=\numberballmDarkTeal]
    %         \item \textbf{\uline{N. Multiplicadores}}
    %               \begin{center}
    %                   \textcolor{blue}{prefijo\subin{1}} \textcolor{blue}{óxido} + de + \textcolor{red}{prefijo\subin{2}} \textcolor{red}{nombre del elemento}
    %               \end{center}

    %               \begin{enumerate}[label=\textcolor{red}{$\bullet$}]
    %                   \item \textcolor{blue}{Prefijo\subin{1}}: Indica el número de oxígenos.
    %                   \item \textcolor{red}{Prefijo\subin{2}}: Indica el número de elementos distintos al oxígeno.
    %               \end{enumerate}
    %     \end{enumerate}
    % \end{alertblock}
\end{frame}

\begin{frame}[t]{Óxidos}
    \begin{alertblock}{\Large Nomenclatura}
        \vspace*{2mm}
        \begin{enumerate}[label=\numberballmDarkTeal]
            \item \textbf{\uline{Nomenclarura de Multiplicadores}}
                  \begin{center}
                      \textcolor{blue}{prefijo\subin{1}} \textcolor{blue}{óxido} + de + \textcolor{red}{prefijo\subin{2}} \textcolor{red}{nombre del elemento}
                  \end{center}

                  \begin{enumerate}[label=\textcolor{red}{$\bullet$}]
                      \item \textcolor{blue}{Prefijo\subin{1}}: Indica el número de oxígenos.
                      \item \textcolor{red}{Prefijo\subin{2}}: Indica el número de elementos distintos al oxígeno.
                  \end{enumerate}
        \end{enumerate}
        \textbf{\uline{Ejemplos}}:\nline

        \ballnumber{1} : \ch{Na2O}: \textcolor{blue}{monóxido de disodio} \ch{->} \textcolor{red}{óxido de disodio}\nline

        \ballnumber{2} : \ch{Ni2O3}: \textcolor{blue}{trióxido de diníquel}
    \end{alertblock}
\end{frame}

\begin{frame}[t]{\worldflag[width=3.9mm,stretch=1]{GB} Oxides}
    \begin{alertblock}{\Large Nomenclature}
        \vspace*{2mm}
        \begin{enumerate}[label=\numberballmDarkTeal]
            \item \textbf{\uline{Nomenclature with multipliers}}
                  \begin{center}
                    \textcolor{red}{prefix\subin{2}} \textcolor{red}{name of the element} 
                    \textcolor{blue}{prefix\subin{1}} \textcolor{blue}{oxide} 
                  \end{center}

                  \begin{enumerate}[label=\textcolor{red}{$\bullet$}]
                      \item \textcolor{blue}{Prefix\subin{1}}: Number of oxygen atoms.
                      \item \textcolor{red}{Prefix\subin{2}}: number of atoms of the element other than oxygen.
                  \end{enumerate}
        \end{enumerate}
        \textbf{\uline{Examples}}:\nline

        \ballnumber{1} : \ch{Na2O}: \textcolor{blue}{disodium monoxide} \ch{->} \textcolor{red}{disodium oxide}\nline

        \ballnumber{2} : \ch{Ni2O3}: \textcolor{blue}{diniquel trioxide}
    \end{alertblock}
\end{frame}

% \begin{frame}[t]{Óxidos}

%     \begin{alertblock}{\Large Nomenclatura}
%         \begin{enumerate}[label=\numberballmDarkTeal]
%             \item \textbf{\uline{Nomenclatura Multiplicadores}}\nline
%         \end{enumerate}
%         \textbf{\uline{Ejemplos}}:\nline

%         \ballnumber{1} : \ch{Na2O}: \textcolor{blue}{monóxido de disodio} \ch{->} \textcolor{red}{óxido de disodio}\nline

%         \ballnumber{2} : \ch{Ni2O3}: \textcolor{blue}{trióxido de diníquel}

%     \end{alertblock}

% \end{frame}

\begin{frame}[t]{Óxidos}
    \begin{alertblock}{\Large Nomenclatura}

        \begin{enumerate}[start=2,label=\numberballmDarkTeal]
            \item \textbf{\uline{Nomenclatura con Estados de Oxidación}}:
                  \begin{center}
                      \textcolor{blue}{óxido} + de + \textcolor{red}{nombre del elemento} + $\Bigl($ \textcolor{green!50!black}{estado de oxidación del elemento en números romanos}$\Bigr)$
                  \end{center}
                  Cuando el elemento tiene tan solo un único estado de oxidación, el estado de oxidación puede omitirse.
        \end{enumerate}

        \textbf{Ejemplos}:

        \ballnumber{1} : \ch{NiO}: \textcolor{blue}{óxido de níquel (II)}

        \ballnumber{2} : \ch{Na2O}: \textcolor{blue}{óxido de sodio (I)} \ch{->} \textcolor{red}{óxido de sodio}, omitimos (I).
    \end{alertblock}
\end{frame}

\begin{frame}[t]{\worldflag[width=3.9mm,stretch=1]{GB} Oxides}
    \begin{alertblock}{\Large Nomenclature} 

        \begin{enumerate}[start=2,label=\numberballmDarkTeal]
            \item \textbf{\uline{Oxidation State Nomenclature}}:
                  \begin{center}
                    \textcolor{red}{name of the element} $\Bigl($\textcolor{green!50!black}{oxidation state of the element in Roman numerals}$\Bigr)$ \textcolor{blue}{oxide} 
                  \end{center}
                  When the element has only one oxidation state, the oxidation state can be omitted
        \end{enumerate}

        \textbf{Examples}:

        \ballnumber{1} : \ch{NiO}: \textcolor{blue}{Niquel oxide (II)}

        \ballnumber{2} : \ch{Na2O}: \textcolor{blue}{sodium oxide (I)} \ch{->} \textcolor{red}{sodium oxide}, we omit the (I).
    \end{alertblock}
\end{frame}

\begin{frame}[t]{Óxidos}
    \begin{exampleblock}{\Large Formulación}
        \begin{enumerate}[label=\numberballmDarkTeal]
            \item Escribimos el oxígeno a la derecha.
            \item Escribimos el elemento a la izquierda.
            \item Ponemos encima de cada elemento su correspondiente estado de oxidación.
            \item Intercambiamos los estados de oxidación como subíndices en valor absoluto.
            \item Simplificamos la fórmula si es posible.
        \end{enumerate}
        \textbf{\uline{Ejemplos}:}\nline
        \ballnumber{1} : Óxido de sodio \ch{->}  \eo{blue}{+1}{Na} \eo{red}{-2}{O}  \ch{->} \ch{Na2O}

        \ballnumber{2} : Óxido de calcio \ch{->} \eo{blue}{+2}{Ca} \eo{red}{-2}{O}  \ch{->} \ch{Ca2O2} \ch{->} \ch{CaO}
    \end{exampleblock}
\end{frame}



\begin{frame}[t]{Oxígeno con halógenos}\vspace{10pt}
    \begin{block}{\Large Características}
        \begin{enumerate}[label=\numberballmDarkTeal]
            \item Formados por \ch{O} y un elemento del grupo 17 (\alert{halógenos}).
            \item Fórmula General: \ch{O_{m}X_{n}}
        \end{enumerate}
    \end{block}
    \begin{alertblock}{\Large Nomenclatura}
        \begin{enumerate}[label=\itembolascolor{\arabic*}{mDarkTeal}]
            \item \textbf{\uline{Nomenclatura con multiplicadores.}}
                  \begin{center}
                      \textcolor{blue}{prefijo\subin{1}} \textcolor{azulF}{raíz del nombre del halógeno} terminado en \textcolor{green!50!black}{uro} + de + \textcolor{red}{prefijo\subin{2}} \textcolor{red}{oxígeno}
                  \end{center}

                  \begin{enumerate}[label=\textcolor{red}{$\bullet$}]
                      \item \textcolor{blue}{Prefijo\subin{1}}: Indica el número de halógenos.
                      \item \rojo{Prefijo\subin{2}}: Indica el número de oxígenos.
                  \end{enumerate}

                  \textbf{\uline{Ejemplos}}: \ballnumber{1} : \ch{OCl2}: dicloruro de oxígeno.
                  \ballnumber{2} : \ch{O3Br2} dibromuro de trioxígeno
        \end{enumerate}
    \end{alertblock}
\end{frame}


\begin{frame}[t]{\worldflag[width=3.9mm,stretch=1]{GB} Oxygen with halogens}\vspace{10pt}
    
    \begin{alertblock}{\Large Nomenclature}
        \begin{enumerate}[label=\itembolascolor{\arabic*}{mDarkTeal}]
            \item \textbf{\uline{Nomenclature with multipliers}}
                  \begin{center}
                      \textcolor{red}{prefix\subin{2}} \textcolor{red}{oxygen} + \textcolor{blue}{prefix\subin{1}}  \textcolor{azulF}{name of the halogen ending in}  \textcolor{green!50!black}{ide}   
                  \end{center}

                  \begin{enumerate}[label=\textcolor{red}{$\bullet$}]
                      \item \textcolor{blue}{Prefijo\subin{1}}: number of halogen atoms.
                      \item \rojo{Prefijo\subin{2}}: number of oxygen atoms.
                  \end{enumerate}

                  \textbf{\uline{Examples}}: \ballnumber{1} : \ch{OCl2}: oxygen dichloride.
                  \ballnumber{2} : \ch{O3Br2} trioxygen dibromide
        \end{enumerate}
    \end{alertblock}
\end{frame}




\begin{frame}[t]{Oxígeno con halógenos}
    \begin{exampleblock}{\Large Formulación}\vspace*{5pt}
        \begin{enumerate}[label=\itembolascolor{\arabic*}{mDarkTeal}]
            \item Escribimos el oxígeno a la izquierda
            \item Escribimos el halógeno a la derecha.
            \item Ponemos encima de cada elemento su correspondiente estado de oxidación.
            \item Intercambiamos los estados de oxidación como subíndices en valor absoluto.
            \item Simplificamos la fórmula si es posible.
        \end{enumerate}
    \end{exampleblock}
    \textbf{\uline{Ejemplos}}

    \ballnumber{1} : Dibromuro de oxígeno: \eob{-2}{O} \eor{+1}{Br}  \ch{-> OBr2}

    \ballnumber{2} : \ch{O7Br2}: \eor{-2}{O}\eob{+7}{Br} \ch{-> O7Br2}
\end{frame}


\subsection{Compuestos con hidrógeno}

\begin{frame}[c]{Compuestos binarios con hidrógeno: hidruros}
    \begin{block}{\Large  Características}\vspace*{5pt}
        \begin{enumerate}[label=\itembolascolor{\arabic*}{mDarkTeal}]
            \item Son compuestos que contienen hidrógeno y un metal o un no metal a \rojo{excepción de los elementos del grupo 16 y 17.}

            \item En este tipo de compuestos el hidrógeno tiene un \azul{estado de oxidación ($-1$)}.

            \item Fórmula general: \ch{M_xH_y}
        \end{enumerate}
    \end{block}
\end{frame}

\begin{frame}[t]{Hidruros metálicos / no metálicos (excepto 16 y 17)}
    \begin{alertblock}{\Large Nomenclatura}\vspace*{5pt}
        \begin{enumerate}[label=\itembolascolor{\arabic*}{mDarkTeal}]
            \item \textbf{\uline{Nomenclatura multiplicadores}}:
                  \begin{center}
                      \textcolor{blue}{prefijo\subin{1}} \textcolor{blue}{hidruro}  + de + \textcolor{red}{nombre del elemento}
                  \end{center}

                  \begin{enumerate}[label=\textcolor{red}{$\bullet$}]
                      \item \azul{Prefijo\subin{1}}: Indica el número de hidrógenos.
                    %   \item \rojo{Prefijo\subin{2}}: Indica el número de elementos distintos al hidrógeno.
                  \end{enumerate}

                  \textbf{\uline{Ejemplos}}:\nline

                  \ballnumber{1} : Trihidruro de aluminio: \ch{AlH3}

                  \ballnumber{2} : Dihidruro de hierro: \ch{FeH2}
        \end{enumerate}
    \end{alertblock}
\end{frame}

\begin{frame}[t]{\worldflag[width=3.9mm,stretch=1]{GB} Metal hydrides / non metal hydrides (except  elements of the grupos 16 y 17)}
    \begin{alertblock}{\Large Nomenclature}\vspace*{5pt}
        \begin{enumerate}[label=\itembolascolor{\arabic*}{mDarkTeal}]
            \item \textbf{\uline{Nomenclature with multipliers}}:
                  \begin{center}
                    \textcolor{red}{prefix\subin{2} name of the element} \textcolor{blue}{prefix\subin{1}} \textcolor{blue}{hydride}

                      \textcolor{blue}{prefix\subin{1}} \textcolor{blue}{hidruro}  + de + \textcolor{red}{prefix\subin{2} nombre del elemento}
                  \end{center}

                  \begin{enumerate}[label=\textcolor{red}{$\bullet$}]
                      \item \azul{Prefijo\subin{1}}: number of hydrogen atoms.
                      \item \rojo{Prefijo\subin{2}}: number of atoms of the element other than oxygen.
                  \end{enumerate}

                  \textbf{\uline{Ejemplos}}:\nline

                  \ballnumber{1} : Aluminium trihydride: \ch{AlH3}

                  \ballnumber{2} : Iron dihydride: \ch{FeH2}
        \end{enumerate}
    \end{alertblock}
\end{frame}

\begin{frame}[t]{Hidros metálicos / no metálicos (excepto 16 y 17)}
    \begin{alertblock}{\Large
            Nomenclatura}\vspace*{5pt}
        \begin{enumerate}[label=\itembolascolor{\arabic*}{mDarkTeal}]
            \addtocounter{enumi}{1}
            \item \textbf{\uline{Nomenclatura estados de oxidación}}
                  \begin{center}
                      \textcolor{blue}{hidruro} + de + \textcolor{red}{nombre del elemento} + $\Bigl($\textcolor{green!50!black}{estado de oxidación del elemento en números romanos}$\Bigr)$
                  \end{center}
        \end{enumerate}


        \textbf{\uline{Ejemplos}}:

        \ballnumber{1} : Hidruro de niquel (II): \ch{NiH2}.

        \ballnumber{2} : Hidruro de plomo (IV): \ch{PbH4}

    \end{alertblock}
\end{frame}

\begin{frame}[t]{\worldflag[width=3.9mm,stretch=1]{GB} Metal hydrides / non metal hydrides (except  elements of the grupos 16 y 17)}
    \begin{alertblock}{\Large Nomenclature} 

        \begin{enumerate}[start=2,label=\numberballmDarkTeal]
            \item \textbf{\uline{Oxidation State Nomenclature}}:
                  \begin{center}
                    \textcolor{red}{name of the element} $\Bigl($\textcolor{green!50!black}{oxidation state of the element in Roman numerals}$\Bigr)$ \textcolor{blue}{hydride} 
                  \end{center}
                  When the element has only one oxidation state, the oxidation state can be omitted
        \end{enumerate}

        \textbf{Examples}:

        \ballnumber{1} : \ch{NiO}: \textcolor{blue}{Niquel (II) hydride (II)}

        \ballnumber{2} : \ch{Na2O}: \textcolor{blue}{sodium (I) hydride} \ch{->} \textcolor{red}{sodium oxide}, we omit the (I).
    \end{alertblock}
\end{frame}



\begin{frame}[t]{Hidros metálicos / no metálicos (excepto 16 y 17)}
    \begin{exampleblock}{\Large{Formulación}}\vspace*{5pt}
        \begin{enumerate}[label=\itembolascolor{\arabic*}{mDarkTeal}]
            \item Ponemos el metal/no metal a la izquierda.
            \item Ponemos el hidrógeno a la derecha.
            \item Ponemos encima de cada elementos los correspondientes estados de oxidación.
            \item Intercambiamos los estados de oxidación en valor absoluto.
            \item Simplificamos la fórmula si es necesario.
        \end{enumerate}

        \textbf{\uline{Ejemplos:}}

        \ballnumber{1} : Hidruro de sodio \ch{->} \eo{blue}{+1}{Na} \eor{-1}{H}  \ch{->} \ch{NaH}

        \ballnumber{2} : Hidruro de calcio \ch{->} \eob{+2}{Ca} \eor{-1}{H}  \ch{->} \ch{CaH2}
    \end{exampleblock}
\end{frame}



\begin{frame}[t]{Hidruros con elementos de los grupos 16 y 17}
    \begin{block}{\Large Características}\vspace*{5pt}
        \begin{enumerate}[label=\itembolascolor{\arabic*}{mDarkTeal}]
            \item En este caso, tendremos compuestos que contienen hidrógeno y elementos de los grupos 16 y 17.
            \item El hidrógeno actúa con estado de oxidación (\rojo{$+1$}). Normalmente los elementos de los grupos 16 y 17 actúan con el \azul{estado de oxidación negativo más bajo}, \rojo{\(-2\)} para los elementos del grupo 16 y \rojo{\(-1\)} para los elementos del grupo 17.
        \end{enumerate}
    \end{block}
\end{frame}

\begin{frame}[t]{Hidruros con elementos de los grupos 16 y 17}
    \begin{alertblock}{\Large Nomenclatura}\vspace*{5pt}
        \begin{enumerate}[label=\itembolascolor{\arabic*}{mDarkTeal}]
            \item \textbf{\uline{Nomenclatura de multiplicadores}}:
                  \begin{center}
                      \textcolor{blue}{prefijo\subin{1}} \textcolor{blue}{raiz del elemento} + \textcolor{mygreen}{uro} + de + \textcolor{red}{prefijo\subin{2}} \textcolor{red}{hidrógeno}
                  \end{center}
                  \begin{enumerate}[label=\textcolor{red}{$\bullet$}]
                      \item \azul{Prefijo\subin{1}}: Indica el número de elementos.
                      \item \rojo{Prefijo\subin{2}}: Indica el número de hidrógenos.
                  \end{enumerate}

                  \textbf{\uline{Ejemplos}}:\nline

                  \ballnumber{1} : \ch{HCl}: cloruro de hidrógeno

                  \ballnumber{2} : \ch{H2Se}: selenuro de dihidrógeno
        \end{enumerate}
    \end{alertblock}
    \textbf{\rojo{\uline{Nota importante}}}: No se suele nombrar con estados de oxidación porque ya conocemos ambos estados de oxidación. Por eso en muchos casos se omiten los prefijos numerales.
\end{frame}

\begin{frame}[t]{\worldflag[width=3.9mm,stretch=1]{GB} Hydrides with elements of the grupos 16 and 17}\vspace{10pt}
    
    \begin{alertblock}{\Large Nomenclature}
        \begin{enumerate}[label=\itembolascolor{\arabic*}{mDarkTeal}]
            \item \textbf{\uline{Nomenclature with multipliers}}
                  \begin{center}

                      \textcolor{red}{prefix\subin{2}} \textcolor{red}{hydrogen} + \textcolor{blue}{prefix\subin{1}}  \textcolor{azulF}{name of the element  ending in}  \textcolor{green!50!black}{ide}   
                  \end{center}

                  \begin{enumerate}[label=\textcolor{red}{$\bullet$}]
                      \item \textcolor{blue}{Prefix\subin{1}}: number of element (non hydrogen) atoms.
                      \item \rojo{Prefix\subin{2}}: number of hydrogen atoms.
                  \end{enumerate}

                  \textbf{\uline{Examples}}:

                  \ballnumber{1} : \ch{HCl}: hydrogen chloride

                  \ballnumber{2} : \ch{H2Se}: dihydrogen selenide
        \end{enumerate}
    \end{alertblock}
\end{frame}

\begin{frame}[t]{Hidruros con elementos de los grupos 16 y 17}
    \begin{exampleblock}{\Large Formulación}\vspace*{5pt}
        \begin{enumerate}[label=\itembolascolor{\arabic*}{mDarkTeal}]
            \item Ponemos el hidrógeno a la izquierda.
            \item Ponemos el elemento del grupo 16 o 17 a la derecha.
            \item Ponemos encima de cada elementos los correspondientes estados de oxidación.
            \item Intercambiamos los estados de oxidación en valor absoluto.
            \item Simplificamos la fórmula si es necesario.
        \end{enumerate}
        \textbf{\uline{Ejemplos}}:

        \ballnumber{1} : Fluoruro de hidrógeno: \ch{HF}

        \ballnumber{2} : Sulfuro de dihidrógeno: \ch{H2S}
    \end{exampleblock}
\end{frame}



\begin{frame}[t]{Hidruros con elementos de los grupos 16 y 17: En disolución acuosa}
    \begin{block}{\Large Ácidos hídricos}\vspace*{5pt}

        \textbf{\coloruwave{red}{En disolución acuosa}}: Los hidruros de los elementos de los grupos 16 y 17 se nombran como ácidos siguiendo la siguiente regla:
        \begin{center}
            \textcolor{blue}{ácido} + \textcolor{red}{raíz del elemento} + \textcolor{green!50!black}{hídrico}
        \end{center}

        \textbf{\uline{Ejemplos}}:

        \ballnumber{1} : \ch{HCl}: ácido clorhídrico

        \ballnumber{2} : \ch{H2S}: ácido sulfhídrico
    \end{block}
\end{frame}

\begin{frame}[t]{\worldflag[width=3.9mm,stretch=1]{GB} Hydrides with elements of the grupos 16 and 17 in water}\vspace{10pt}
    
    \begin{alertblock}{\Large Nomenclature}\vspace*{6pt}
            \textbf{\coloruwave{red}{Aqueous solution}}: the hydrides of the elements in groups 16 and 17 are named following the following rule:
                  \begin{center}
                    \textcolor{green!50!black}{hydro} \textcolor{blue}{prefix\subin{1}}  \textcolor{azulF}{name of the element ending in idic} \textcolor{green!50!black}{ide}  acid 
                  \end{center}


                  \textbf{\uline{Examples}}:

                  \ballnumber{1} : \ch{HCl}: hydrochloridic acid

                  \ballnumber{2} : \ch{H2S}: hydrosulfidic acid
    \end{alertblock}
\end{frame}

\begin{frame}[t]{Hidruos no metálicos}
    \begin{block}{\Large Características}\vspace*{5pt}
        \begin{enumerate}
            \item Combinaciones de hidrógeno con elementos no metálicos de los grupos 13, 14 y 15.
            \item Actuará con \rojo{($-1$)} con los elementos de los grupos 13, 14, y 15.
            \item Nomenclatura de multiplicadores como hidruros metálicos.
            \item Nomenclatura tradicional.
        \end{enumerate}
        Ambas nomenclaturas están resumidas en la siguiente tabla.
    \end{block}
\end{frame}

\begin{frame}
    \begin{tblr}[caption={Hidruros}]
        {colspec={Q[c,m] Q[c,m] Q[c,m]},
            row{1} = {bg=myOrange3!90!white, fg = white, font=\sffamily}
        }
        \hline
        Compuesto & N. Multiplicadores \TblrNote{a} & N. Tradicional \TblrNote{b} \\
        \hline
        \ch{BH3}  & trihidruro de boro              & borano                      \\
        \ch{CH4}  & tetrahidruro de carbono         & metano                      \\
        \ch{SiH4} & tetrahidruro de sicilio         & silano                      \\
        \ch{NH3}  & trihiruro de nitrógeno          & amoniaco o azano            \\
        \ch{PH3}  & trihidruro de fósforo           & fosfano                     \\
        \ch{AsH3} & trihidruro de arsénico          & arsano                      \\
        \ch{SbH3} & trihidruro de antimonio         & estibano                    \\
        \hline
    \end{tblr}
\end{frame}



\subsection{Combinaciones Metal + no Metal}

\begin{frame}[c]{Metal + no Metal}
    \begin{block}{\Large Características}
        \begin{enumerate}[label=\itembolascolor{\arabic*}{mDarkTeal}]
            \item Metal + no metal
            \item El no metal suele actuar con el estado de oxidación más bajo, es decir el más negativo. Se pondrá a la derecha.
            \item El metal se pondrá a la izquierda y tendrá el estado de oxidación positivo.
        \end{enumerate}
    \end{block}
\end{frame}

\begin{frame}[t]{Metal + no Metal}
    \begin{alertblock}{\Large Nomenclatura}\vspace*{5pt}

        \begin{enumerate}[label=\itembolascolor{\arabic*}{mDarkTeal}]
            \item \textbf{\uline{Nomenclatura de multiplicadores}}:
                  \begin{center}
                      \textcolor{blue}{prefijo\subin{1}} \textcolor{blue}{raiz del no metal} \textcolor{mygreen}{uro} + de + \textcolor{red}{prefijo\subin{2}} \textcolor{red}{nombre del metal}
                  \end{center}
                  \begin{enumerate}[label=\textcolor{red}{$\bullet$}]
                      \item \azul{Prefijo\subin{1}}: Indica el número de elementos no metálicos.
                      \item \rojo{Prefijo\subin{2}}: Indica el número de elementos metálicos.
                  \end{enumerate}
        \end{enumerate}
        \textbf{\uline{Ejemplos}}

        \ballnumber{1} : \ch{Al4C3}: tricarburo de tetraaluminio

        \ballnumber{2} : \ch{SnTe2}: diteluro de estaño
    \end{alertblock}
\end{frame}

\begin{frame}[t]{Metal + no Metal}
    \begin{alertblock}{\Large Nomenclatura}\vspace*{5pt}

        \begin{enumerate}[label=\itembolascolor{\arabic*}{mDarkTeal}]
            \addtocounter{enumi}{1}
            \item \textbf{\uline{Nomenclatura de Estados de Oxidación}}:
                  \begin{center}
                      \textcolor{blue}{raiz del no metal} \textcolor{mygreen}{uro}  + \textcolor{red}{nombre del metal} + $\Bigl($\textcolor{green!50!black}{estado de oxidación del metal en números romanos}$\Bigr)$
                  \end{center}

        \end{enumerate}
        \textbf{\uline{Ejemplos}}

        \ballnumber{1} : \ch{Al4C3}: carburo de aluminio (III) \ch{->} carburo de alumnio

        \ballnumber{2} : \ch{SnTe2}: teluro  de estaño (IV)
    \end{alertblock}
\end{frame}


\begin{frame}[t]{Metal + no Metal}
    \begin{exampleblock}{\Large Formulación}
        \begin{enumerate}[label=\itembolascolor{\arabic*}{mDarkTeal}]
            \item Ponemos el metal a la izquierda.
            \item Ponemos el no metal a la derecha.
            \item Ponemos encima de cada elementos los correspondientes estados de oxidación.
            \item Intercambiamos los estados de oxidación en valor absoluto.
            \item Simplificamos la fórmula si es necesario.
        \end{enumerate}
    \end{exampleblock}
    \textbf{\uline{Ejemplos}}\nline{}
    \ballnumber{1} : Cloruro de sodio \ch{->} \eob{+1}{Na} \eor{-1}{Cl} \ch{->} \ch{NaCl}

    \ballnumber{2} : \ch{->} \eob{+2}{Sn} \eor{-2}{S}  \ch{->} \ch{Sn2S2} \ch{->} \ch{SnS}
\end{frame}



\subsection{Peróxidos}

\begin{frame}[t]{Peróxidos}
    \begin{block}{\Large Características}\vspace*{5pt}
        \begin{enumerate}[label=\itembolascolor{\arabic*}{mDarkTeal}]
            \item Formados por un ión peróxido (\ch{O2^{2-}}) y un catión (un elemento normalmente metálico)
            \item En el ión peróxido el oxígeno tiene estado de oxidación ($-1$)
            \item Esto es útil porque si no se está seguro si es un peróxido o un óxido contamos las cargas y si el oxígeno está con ($-1$) es un peróxido y si está con ($-2$) es un óxido.
            \item Si se simplifica tiene que ser de manera que \blueremark{no} queden menos de dos átomos de oxígeno (un grupo peróxido completo).
        \end{enumerate}

    \end{block}
\end{frame}


\begin{frame}[t]{Peróxidos}
    \begin{alertblock}{\Large Nomenclatura}\vspace*{5pt}
        \begin{enumerate}[label=\itembolascolor{\arabic*}{mDarkTeal}]
            \item \textbf{\uline{Nomenclatura multiplicadores}}: Se nombran siguiendo las reglas de los óxidos.

                  \textbf{\uline{Ejemplos}}:\nline{}

                  \ballnumber{1} : \ch{Li2O2}: dióxido de dilitio
            \item \textbf{\uline{Nomenclatura de peróxidos}}:
                  \begin{center}
                      \textcolor{blue}{peróxido} + de + \textcolor{red}{nombre del elemento}
                  \end{center}

                  \textbf{\uline{Ejemplos}}:\nline{}

                  \ballnumber{1} : \ch{Na2O2}: peróxido de sodio
        \end{enumerate}

    \end{alertblock}
\end{frame}


\begin{frame}[t]{Peróxidos}
    \begin{exampleblock}{\Large Formulación}\vspace*{5pt}
        \begin{enumerate}[label=\itembolascolor{\arabic*}{mDarkTeal}]
            \item Ponemos el metal a la izquierda.
            \item Ponemos el grupo peróxido \ch{O2^{-2}} a la derecha.
            \item Ponemos encima de cada elementos los correspondientes estados de oxidación.
            \item Intercambiamos los estados de oxidación en valor absoluto.
            \item Simplificamos la fórmula si es necesario.
        \end{enumerate}

        \textbf{\uline{Ejemplos}}

        \ballnumber{1} : Peróxido de sodio \ch{->} \eob{+1}{Na} \eor{-2}{O2}  \ch{->} \ch{Na2O2}

        \ballnumber{2}: Peróxido de aluminio: \ch{->} \eob{+3}{Al} \eor{-2}{O2}  \ch{->} \ch{Al2(O2)3}


    \end{exampleblock}
\end{frame}



\subsection{Hidróxidos}

\begin{frame}[t]{Hidróxidos}
    \begin{block}{\Large Características}\vspace*{5pt}
        \begin{enumerate}[label=\itembolascolor{\arabic*}{mDarkTeal}]
            \item Compuestos ternarios con tres elementos.
            \item Formados por un metal (o catión) y el anión hidróxido  \ch{OH-}
            \item El grupo  \ch{OH-} lo debemos considerar como un grupo indivisible teniendo como estado de oxidación $-1$.
            \item La fórmula general para los hidróxidos será \ch{M(OH)_n}. Siendo M, el metal y $n$ estado de oxidación intercambiado por el metal.
        \end{enumerate}
    \end{block}
\end{frame}

\begin{frame}[t]{Hidróxidos}
    \begin{alertblock}{\Large Nomenclatura}\vspace*{5pt}
        \begin{enumerate}[label=\itembolascolor{\arabic*}{mDarkTeal}]
            \item \textbf{\uline{Nomenclatura multiplicadores}}:
                  \begin{center}
                      \textcolor{blue}{prefijo\subin{1} hidróxido} + de + \textcolor{red}{nombre del elemento distinto al OH}
                  \end{center}

                  \begin{enumerate}[label=\textcolor{red}{$\bullet$}]
                      \item \azul{Prefijo\subin{1}}: Indica el número de grupos \ch{OH}.
                    %   \item \rojo{Prefijo\subin{2}}: Indica el número de elementos metálicos.
                  \end{enumerate}
        \end{enumerate}

        \textbf{\uline{Ejemplos}}:

        \ballnumber{1} : \ch{LiOH}: hidróxido de litio

        \ballnumber{2} : \ch{Mg(OH)2}: dihidróxido de magnesio
    \end{alertblock}

\end{frame}

\begin{frame}[t]{Hidróxidos}
    \begin{alertblock}{\Large Nomenclatura}\vspace*{5pt}
        \begin{enumerate}[label=\itembolascolor{\arabic*}{mDarkTeal}]
            \addtocounter{enumi}{1}
            \item \textbf{\uline{Nomenclatura Estados de oxidación}}:
                  \begin{center}
                      \textcolor{blue}{hidróxido} + de + \textcolor{red}{nombre del elemento} + $\Bigl($\textcolor{green!50!black}{estado de oxidación en números romanos del elemento}$\Bigr)$
                  \end{center}
        \end{enumerate}
        \textbf{\uline{Ejemplos}}:

        \ballnumber{1} : \ch{Fe(OH)2}: \azul{Hidróxido de hierro (II)}

        \ballnumber{2} : \ch{CuOH}: \azul{Hidróxido de cobre (I)}
    \end{alertblock}
\end{frame}

\begin{frame}[c]{Hidróxidos}
    \begin{exampleblock}{\Large Formulación}\vspace*{5pt}
        \begin{enumerate}[label=\itembolascolor{\arabic*}{mDarkTeal}]
            \item Ponemos el metal a la izquierda.
            \item Ponemos el grupo hidróxido \ch{OH-} a la derecha.
            \item Ponemos encima de cada elementos los correspondientes estados de oxidación.
            \item Intercambiamos los estados de oxidación en valor absoluto.
            \item Ponemos entre paréntesis el grupo \ch{OH} y a continuación el estado de oxidación intercambiado. El paréntesis se omite si el estado de oxidación intercambiando es $1$.
            \item Simplificamos la fórmula si es necesario.
        \end{enumerate}
    \end{exampleblock}
\end{frame}

\begin{frame}[c]{Hidróxidos}
    \begin{exampleblock}{\Large Formulación}\vspace*{5pt}

        \textbf{\uline{Ejemplos}}:

        \ballnumber{1}: Hidróxido de sodio: \eob{+1}{Na} \eor{-1}{OH}  \ch{->} \ch{NaOH}

        \ballnumber{2}: Hidróxido de calcio: \eob{+2}{Ca} \eor{-1}{OH}  \ch{->} \ch{Ca(OH)2}
    \end{exampleblock}
\end{frame}

\section{Compuestos terciarios}
\subsection{Oxoácidos}

\begin{frame}[c]{Oxoácidos}
    \begin{block}{\Large Características}
        \begin{enumerate}[label=\itembolascolor{\arabic*}{mDarkTeal}]
            \item Son compuestos ternarios,son generalmente el hidrógeno, un no metal (o un metal de transición como el Cr o el Mn) y oxígeno.
            \item Tienen como fórmula general \ch{H_iX_jO_k}, siendo X el no metal o metal de transición.
            \item El hidrógeno siempre tendrá un estado de oxidación ($+1$) y el oxígeno ($-2$)
        \end{enumerate}

    \end{block}
\end{frame}

\begin{frame}[c]{Oxoácidos}
    \begin{exampleblock}{\Large Formulación}
        \begin{enumerate}[label=\itembolascolor{\arabic*}{mDarkTeal}]
            \item Ponemos el hidrógeno a la izquierda.
            \item Ponemos el elemento no metálico o metal de transición en la parte central.
            \item Ponemos el oxígeno en la parte derecha.
        \end{enumerate}
    \end{exampleblock}
\end{frame}

\begin{frame}[t]{Oxoácidos}
    \begin{exampleblock}{\Large Formulación}
        \begin{enumerate}[label=\itembolascolor{\arabic*}{mDarkTeal}]
            \addtocounter{enumi}{3} % para continuar en la numeración de donde quiero
            \item Los oxoácidos provinene de la reacción entre un óxido y el agua como por ejemplo:
                  \begin{center}
                      \ch{SO + H2O -> H2SO2}\\
                      \ch{SO2 + H2O -> H2SO3}\\
                      \ch{SO3 + H2O -> H2SO4}\\
                      \ch{N2O5 + H2O -> 2 HNO3}\\
                      \ch{OCl2 + H2O -> H2Cl2O2 -> 2 HClO}\\
                      \ch{O3Cl2 + H2O -> H2Cl2O4 -> 2 HClO2}\\
                      \ch{O5Cl2 + H2O -> H2Cl2O6 -> 2 HClO3}\\
                      \ch{O7Cl2 + H2O -> H2Cl2O8 -> 2 HClO4}
                  \end{center}
        \end{enumerate}
    \end{exampleblock}
\end{frame}

\begin{frame}[t]{Oxoácidos}
    \begin{exampleblock}{\Large Formulación}
        \vspace*{6pt}
        Podemos estimar la fórmula general de cada oxoácido según se muestra en la siguiente tabla, (siendo X el elemento no metálico o metal de transición):
        \begin{center}
            \begin{figure}[H]
                \centering
                \SetTblrInner{rowsep=4pt}
                \begin{tblr}[note{a}= {Estado de oxidación impar},
                        note{b}= {Estado de oxidación par}]{colspec={Q[c,m,2cm] Q[c,m, 2cm] | Q[c,m, 2cm] Q[c,m, 2cm]},
                        row{1} = {bg=myOrange3!90!white, fg = white, font=\sffamily}
                    }
                    \hline
                    E.O impar \TblrNote{a} & Fórmula General & E.O par \TblrNote{b} & Fórmula General \\
                    \hline
                    +1                     & \ch{HXO}        & +2                   & \ch{H2XO2}      \\
                    +3                     & \ch{HXO2}       & +4                   & \ch{H2XO3}      \\
                    +5                     & \ch{HXO3}       & +6                   & \ch{H2XO4}      \\
                    +7                     & \ch{HXO4}       &                      &                 \\
                    \hline
                \end{tblr}
            \end{figure}
        \end{center}
    \end{exampleblock}
\end{frame}


\begin{frame}[c]{Oxoácidos}
    \begin{exampleblock}{\Large Formulación}
        \vspace*{6pt}
        \textbf{\uline{\textcolor{red}{Nota Importante}}}:
        \begin{enumerate}[label=\itembolascolor{\arabic*}{mDarkTeal}]
            \item Los elementos con estados de oxidación impares, tienen un número impar de átomos de hidrógeno. (\rojo{normalmente 1}).
            \item Los elementos con estado de oxidación pares tienen un número par de átomos de hidrógeno (\rojo{normalmente 2}).
        \end{enumerate}
    \end{exampleblock}
\end{frame}


\begin{frame}[t]{Oxoácidos}
    \examplebeamer{\Large Formulación}
    \begin{enumerate}[label=\numberballmDarkTeal]
        \item El \negrita{\rojo{boro, fósforo, arsénico y silicio}} dan, además de esos compuestos otros ácidos que se forman añadiendo una molécula de agua a los anteriores (al correspondiente ácido formado a partir del óxido y $1$ molécula de agua).
        \item Estos son más estables y se pueden considear los verdaderos ácidos del boro, fósforo, arsénico y silicio.
    \end{enumerate}

    \begin{center}
        \ch{HPO2 + H2O -> H3PO3}\nline{}
        \ch{HPO3 + H2O -> H3PO4}\nline{}
        \ch{HBO2 + H2O -> H3BO3}\nline{}
        \ch{H2SiO3 +  H2O -> H4SiO4}
    \end{center}
\end{frame}

\begin{frame}[t]{Oxoácidos}
    La tabla anterior por tanto la podemos sinterizar en la siguiente:
    \begin{center}
        \begin{figure}[H]
            \centering
            \SetTblrInner{rowsep=4pt}
            \begin{tblr}[note{a} = {Estado de oxidación},
                    note{b} = {Fórmula general}]
                {colspec={Q[c,m] Q[c,m] Q[c,m] Q[c,m] Q[c,m] Q[c,m]},
                    row{1} = {bg=myOrange3!90!white, fg = white, font=\sffamily}
                }
                \hline
                E.O \TblrNote{a} & F. General & Para B, P, As & E.O & F. General & Para Si     \\
                \hline
                +1               & \ch{HXO}   & ---           & +2  & \ch{H2XO2} & ---         \\
                +3               & \ch{HXO2}  & \ch{H3XO3}    & +4  & \ch{H2XO3} & \ch{H4SiO4} \\
                +5               & \ch{HXO3}  & \ch{H3XO4}    & +6  & \ch{H2XO4} & ---         \\
                +7               & \ch{HXO4}  & ---           &     &            &             \\
                \hline
            \end{tblr}
        \end{figure}
    \end{center}
\end{frame}



\begin{frame}[t]{Oxoácidos}
    \alertbeamer{\Large Nomenclatura}
    \begin{enumerate}[label=\numberballmDarkTeal]
        \item \textbf{\uline{Nomcenclatura Tradicional}} Es la más utilizada y admitida por la IUPAC. Se nombran según la siguiente regla general haciendo uso de los prefijos \blueremark{hipo} o \blueremark{per} y los sufijos \blueremark{oso} y \blueremark{ico}.
              \begin{center}
                  \textcolor{blue}{ácido} + \textcolor{mygreen}{raíz del elemento} + \textcolor{red}{oso/ico}\par
                  o también\par
                  \textcolor{blue}{ácido} + \textcolor{red}{hipo/per} + \textcolor{mygreen}{raíz del elemento} + \textcolor{red}{oso/ico}
              \end{center}
    \end{enumerate}
    El uso de estos sufijos y prefijos sirven para indicar el estado de oxidación con el que actúa el elemento central según la siguiente tabla.

\end{frame}

\begin{frame}[t]{Oxoácidos}
    \alertbeamer{\Large Nomenclatura}.
    \begin{enumerate}[label=\numberballmDarkTeal]
        \item\textbf{\uline{Nomcenclatura Tradicional}}
    \end{enumerate}
    \begin{center}
        \begin{figure}[H]
            \centering
            \begin{tblr}{colspec={Q[c,m] Q[c,m] Q[c,m] Q[c,m] Q[c,m]}}
                \hline
                Elementos                & hipo \dots oso & oso & ico & per \dots ico \\
                \hline
                Grupo 14 (C, Si)        & ---            & +2  & +4  & ---           \\
                Grupo 15 (N, P, As, Sb) & +1             & +3  & +5  & ---           \\
                Grupo 16 (S, Se, Te)    & +2             & +4  & +6  & ---           \\
                Grupo 17 (F, Cl, Br, I) & +1             & +3  & +5  & +7            \\
                Mn                      & ---            & +4  & +6  & +7            \\
                Cr, Mo, W               & ---            & --- & +6  & ---           \\
                V                       & ---            & --- & +5  & ---           \\
                \hline
            \end{tblr}
        \end{figure}
    \end{center}

\end{frame}

\begin{frame}[t]{Oxoácidos}
    \alertbeamer{\Large Nomenclatura}\nline{}
    \blockbeamer{\uline{Reglas a tener en cuenta}}
    \begin{enumerate}[label=\numberballmDarkTeal]
        \item Asignamos el sufijo \blueremark{-ico} al estado de oxidación que coincida con el último número del grupo al que pertenezca el elemento.
        \item En el caso del carbono, como pertenece algrupo $14$, le asignaremos el prefijo \blueremark{-ico} cuando el carbono tenga un estado de oxidación de $+4$.
              En el caso del nitrógeno que pertenece al grupo $15$ le asignaremos el prefijo \blueremark{-ico} cuando se encuentre en el estado de oxidación $+5$.
        \item Para los del grupo 16, es decir azufre, selenio \dots se le asignará el sufijo \blueremark{-ico} cuando se encuentre con estado de oxidación $+6$.
    \end{enumerate}
\end{frame}

\begin{frame}[t]{Oxoácidos}
    \alertbeamer{\Large Nomenclatura}\nline{}
    \blockbeamer{\uline{Reglas a tener en cuenta}}
    \begin{enumerate}[label=\numberballmDarkTeal]
        \addtocounter{enumi}{3}
        \item La excepción la tendremos cuando tengamos elementos del grupo 17 (F, Cl, Br, I) en el que en este caso se le asignará el sufijo \blueremark{-ico} cuando se encuentre con estado de oxidación $+5$.
        \item Tomamos como referencia el sufijo \blueremark{-ico} y a partir de este asignaremos el resto de sufijos con el correspondiente estado de oxidación.
        \item Si el estado de oxidación es un número menor al tomado como referencia, se le asigna el sufijo \blueremark{-oso}.
    \end{enumerate}
\end{frame}

\begin{frame}[t]{Oxoácidos}
    \alertbeamer{\Large Nomenclatura}\nline{}
    \blockbeamer{\uline{Reglas a tener en cuenta}}
    \begin{enumerate}[label=\numberballmDarkTeal]
        \addtocounter{enumi}{6}
        \item Si el estado de oxidación son dos  números menor al tomado como referencia, se le asigna el prefijo \blueremark{hipo-} + el sufijo \blueremark{-oso}.
        \item Si el estado de oxidación es un número por encima del estado de oxidación tomado como referencia se le asigna el prefijo \blueremark{per-} y el sufijo \blueremark{-ico}
    \end{enumerate}
\end{frame}

\begin{frame}[t]{Oxoácidos}
    \alertbeamer{\Large Nomenclatura}\nline{}
    \numeromDarkTeal{2} \textbf{\uline{Nomenclatura del Hidrógeno}}:
    Esta nomenclatura es muy recomendada por la IUPAC e indica la composición del ácido Sigue la siguiente regla general:
    \begin{center}
        \textcolor{blue}{prefijo\textsubscript{1}} \textcolor{red}{hidrogeno} $\Bigl($ \textcolor{mygreen}{prefijo\textsubscript{2}} \textcolor{blue}{oxido} \textcolor{myBlue}{prefijo\textsubscript{3}} \textcolor{red}{raiz del nombre del elemento central} \textcolor{mygreen}{ato}$\Bigr)$
    \end{center}
    \begin{enumerate}[label=\textcolor{red}{$\bullet$}]
        \item \azul{Prefijo\textsubscript{1}}: Indica el número de hidrógenos.
        \item \textcolor{mygreen}{Prefijo\textsubscript{2}}: Indica el número de oxígenos.
        \item \textcolor{myBlue}{Prefijo\textsubscript{3}}: Indica el número de átomos del elemento central.
    \end{enumerate}
    Normalmente el número del átomo central es $1$ y por tanto se suele omitir (el prefijo mono se omite si no hay ambigüedad)
\end{frame}

\begin{frame}[t]{Oxoácidos}
    \alertbeamer{\Large Nomenclatura}\nline{}
    \numeromDarkTeal{2} \textbf{\uline{Nomenclatura del Hidrógeno}}:\nline{}
    \textbf{\uline{Ejemplos:}}

    \ballnumber{1} : \ch{H2CO3}: dihidrogeno(trioxidocarbonato)

    \ballnumber{2} : \ch{H2CrO4}: dihidrogeno(tetraoxidocromato)

    \ballnumber{3} : \ch{HNO2}: hidrogeno(dioxidonitrato)

    \textbf{\textcolor{red}{\uline{Nota Importante}}}: Los nombres siempre terminan en \blueremark{-ato}, en este caso no se utiliza -ico, -oso, hipo-, per-. Los términos \blueremark{hidrogeno} y \blueremark{oxido} van sin acento al ser utilizados como prefijos.
\end{frame}


\begin{frame}[t]{Oxoácidos}
    \alertbeamer{\Large Consejos prácticos}\nline{}
    \examplebeamer{\uline{Determinar la fórmula de los oxoácidos a partir del nombre tradicional}}\nline{}
    \blockbeamer{Elementos con estado de oxidación par}. \rojo{Ácido sulfuroso}
    \begin{enumerate}[label=\numberballmDarkTeal]
        \item Primero determinaremos los posibles estados de oxidación del azufre: $+2\ +4\ +6$.
        \item Nos fijamos en el sufijo del elemento central, \blueremark{-oso} indica que es un número menor (ver arriba) al estado de oxidación que tomamos como referencia, en este caso el estado oxidación de referencia sería el $+6$ ya que el azufre se encuentra en el grupo $16$ por lo que el azufre tendrá un estado de oxidación de $+4$.
    \end{enumerate}
\end{frame}

\begin{frame}[t]{Oxoácidos}
    \alertbeamer{\Large Consejos prácticos}\nline{}
    \examplebeamer{\uline{Determinar la fórmula de los oxoácidos a partir del nombre tradicional}}\nline{}
    \blockbeamer{Elementos con estado de oxidación par}. \rojo{Ácido sulfuroso}
    \begin{enumerate}[label=\numberballmDarkTeal]
        \addtocounter{enumi}{2}
        \item Como el estado de oxidación del elemento central es \textit{par} entonces tendremos \textcolor{red}{\uline{2 hidrógenos}} pudiendo escribir: \ch{H2SO_x}.
        \item Para determinar el valor de $x$ se puede hacer de una forma muy sencilla: al estado de oxidación del elemento central se le suma $2$ (\textit{cuando el estado de oxidación es par siempre se suma 2}). En este caso tendríamos que $4 + 2 = 6$.
    \end{enumerate}
\end{frame}

\begin{frame}[t]{Oxoácidos}
    \alertbeamer{\Large Consejos prácticos}\nline{}
    \examplebeamer{\uline{Determinar la fórmula de los oxoácidos a partir del nombre tradicional}}\nline{}
    \blockbeamer{Elementos con estado de oxidación par}. \rojo{Ácido sulfuroso}
    \begin{enumerate}[label=\numberballmDarkTeal]
        \addtocounter{enumi}{4}
        \item Ahora, dividimos entre $2$ el número obtenido en el paso anterior obteniendo que:
              \[\frac{6}{2} = 3\]
              siendo $3$ el valor de $x$.
        \item La formula final quedará entonces como: \ch{H2SO3}
    \end{enumerate}
\end{frame}

\begin{frame}[t]{Oxoácidos}
    \alertbeamer{\Large Consejos prácticos}\nline{}
    \examplebeamer{\uline{Determinar la fórmula de los oxoácidos a partir del nombre tradicional}}\nline{}
    \blockbeamer{Elementos con estado de oxidación impar}. \rojo{Ácido nítrico}
    \begin{enumerate}[label=\numberballmDarkTeal]
        \item Primero determinaremos los posibles estados de oxidación del nitrógeno: $+1\ +3\ +5$.
        \item Nos fijamos en el sufijo del elemento central, \blueremark{-ico} indica que es el número del estado de oxidación que tomamos como referencia, en este caso el estado oxidación de referencia sería el $+5$ ya que el nitrógeno se encuentra en el grupo $15$.
        \item Como el estado de oxidación del elemento central es \textit{impar} entonces tendremos \textcolor{red}{\uline{1 hidrógeno}} pudiendo escribir: \alert{\ch{HNO_x}}.
    \end{enumerate}
\end{frame}

\begin{frame}[t]{Oxoácidos}
    \alertbeamer{\Large Consejos prácticos}\nline{}
    \examplebeamer{\uline{Determinar la fórmula de los oxoácidos a partir del nombre tradicional}}\nline{}
    \blockbeamer{Elementos con estado de oxidación impar}. \rojo{Ácido nítrico}
    \begin{enumerate}[label=\numberballmDarkTeal]
        \addtocounter{enumi}{3}
        \item Para determinar el valor de $x$ se puede hacer de una forma muy sencilla: al estado de oxidación del elemento central se le suma $1$ (\textit{cuando el estado de oxidación es impar siempre se suma 1}). En este caso tendríamos que $5 + 1 = 6$.
        \item Ahora, dividimos entre $2$ el número obtenido en el paso anterior obteniendo que:
              \(\dfrac{6}{2} = 3\) siendo $3$ el valor de $x$.
        \item La formula final quedará entonces como: \ch{HNO3}
    \end{enumerate}
\end{frame}

\begin{frame}[c]{Oxoácidos}
    \alertbeamer{\Large Consejos prácticos}\nline{}
    \examplebeamer{\uline{Determinar la fórmula de los oxoácidos a partir del nombre tradicional}}\nline{}
    \textbf{\textcolor{red}{\uline{Nota Importante}}}: Anteriormente se ha indicado que los elementos B, P, As o Si, forman oxoácidos más estables cuando se añade una molécula más de \ch{H2O} por lo que si aplicamos las reglas anteriores tendremos que sumar una molécula de \ch{H2O} a la fórmula obtenida. Si lo que queremos obtener es el compuesto \blueremark{meta-} entonces si dejaríamos el compuesto tal y como nos sale aplicando el procedimiento descrito.
\end{frame}

\begin{frame}[c]{Oxoácidos}
    \alertbeamer{\Large Consejos prácticos}\nline{}
    \examplebeamer{\uline{Determinar el nombre de los oxoácidos a partir de la fórmula}}. \rojo{\ch{HIO4}}\nline{}
    \begin{enumerate}[label=\numberballmDarkTeal]
        \item Determinar el estado de oxidación del átomo central (I). La suma de todos los estados de oxidación debe ser cero (compuesto neutro).
              \begin{equation*}
                  \underbrace{1\cdot(+1)}_{\text{H}} + \overbrace{4\cdot(-2)}^{\text{O}} + x = 0
              \end{equation*}
              \begin{equation*}
                  x = 6+1 = +7
              \end{equation*}
              Quedando finalmente el estado de oxidación del yodo como ($+7$).
    \end{enumerate}
\end{frame}

\begin{frame}[c]{Oxoácidos}
    \alertbeamer{\Large Consejos prácticos}\nline{}
    \examplebeamer{\uline{Determinar el nombre de los oxoácidos a partir de la fórmula}}. \rojo{\ch{HIO4}}\nline{}
    \begin{enumerate}[label=\numberballmDarkTeal]
        \addtocounter{enumi}{1}
        \item Como el estado de oxidación obtenido ($+7$) es un valor superior al estado de oxidación de referencia que vimos anteriormente para los elementos del grupo 17, entonces tendremos que aplicar los prefijos y sufijos: \blueremark{per-} \dots \blueremark{-ico}.
        \item El nombre por tanto quedaría como: \textcolor{blue}{ácido peryódico}
    \end{enumerate}
\end{frame}


\subsection{Oxosales}

\begin{frame}[t]{Oxosales}
    Son derivados de compuestos ternarios derivados de sustituir todos o algunos de los hidrógenos de los oxoácidos por cationes metálicos como el \ch{Na+} o no metálicos como el ión amonio \ch{NH4+}
    \begin{center}
        \ch{H2SO4 -> 2 H+ + SO4^2-}

        \ch{2 Na+ -> Na2SO4}
    \end{center}

    Para su formación se siguen las mismas reglas que para los ácidos de los que proviene, pero cambiando la terminación \blueremark{-ico} por \blueremark{-ato} y la terminación \blueremark{-oso} por \blueremark{-ito}

\end{frame}

\begin{frame}[t]{Oxosales}
    \examplebeamer{\Large Nomenclatura}

    Utilizaremos como ejemplo \textcolor{blue}{\ch{NaClO4}}
    \begin{enumerate}[label=\numberballmDarkTeal]
        \item Disociamos la sal en sus iones: \ch{Na+} y \ch{ClO4-}
              La carga del anión se puede deducir a partir de la carga del catión, o teniendo en cuenta los hidrógenos que tiene el ácido del que procede.
        \item Formulamos el oxoácido del que procede el oxoanión y lo nombramos:
              \ch{HClO4 -> } ácido perclórico (el cloro actúa con el estado de oxidación mayor ($+7$).
    \end{enumerate}
\end{frame}

\begin{frame}[t]{Oxosales}
    \examplebeamer{\Large Nomenclatura}

    Utilizaremos como ejemplo \textcolor{blue}{\ch{NaClO4}}
    \begin{enumerate}[label=\numberballmDarkTeal]
        \addtocounter{enumi}{2}
        \item Nombramos el oxoanión cambiando el sufijo como hemos dicho anteriormente:
              \ch{ClO4 \mch -> }anión perclorato.
        \item Nombramos la oxosal de derecha a izquierda. Si el metal puede tener otros números de oxidación se indica entre paréntesis el que tiene el compuesto.
              \ch{NaClO4 -> }perclorato de sodio
    \end{enumerate}
\end{frame}

\end{document}